\documentclass[11pt]{article}

% ============================================================
% Milestone 3 -- Final Report Template
% DIP Team Project
% ============================================================
\usepackage[margin=1in]{geometry}
\usepackage[T1]{fontenc}
\usepackage[utf8]{inputenc}
\usepackage{amsmath,amssymb}
\usepackage{graphicx}
\usepackage{booktabs}
\usepackage[round]{natbib}
\usepackage{hyperref}
\hypersetup{colorlinks=true, linkcolor=blue!50!black, citecolor=blue!50!black, urlcolor=blue!50!black}
\usepackage{enumitem}
\usepackage{xcolor}
\usepackage{caption}
\usepackage{subcaption}

\newcommand{\todo}[1]{\textcolor{red}{\textbf{[TODO: #1]}}}

\title{\LARGE \todo{Project Title} \\[4pt]
\large Final Report}
\author{
  \todo{Student A Full Name} \and
  \todo{Student B Full Name} \and
  \todo{Student C Full Name}
}
\date{Milestone 3 --- \todo{Date} \\{} [COURSE CODE] --- Digital Image Processing}

\begin{document}
\maketitle

\begin{abstract}
\todo{150--250 words. Problem, approach, dataset, headline quantitative result vs.
the stated baseline/SOTA reference, and one-sentence takeaway. Write this last.}
\end{abstract}

\section{Introduction}
\label{sec:intro}
\todo{What problem, why it matters, what you did, and a one-paragraph summary of your
findings (reviewers should be able to read the intro and abstract alone and know what
you found).}

\section{Related Work}
\label{sec:related}
\todo{Position your work relative to the references identified at M1 (and any added
since). Group by theme, not by chronology. At least 2 references should be from
2023--2026.}

\section{Method}
\label{sec:method}
\todo{Describe your final approach precisely enough that another team could
reimplement it: architecture, loss function(s), training procedure, hyperparameters.
Use an algorithm block or architecture figure if it helps.}

\begin{figure}[h]
\centering
% \includegraphics[width=0.85\linewidth]{figures/architecture.png}
\caption{\todo{Architecture / pipeline diagram.}}
\label{fig:arch}
\end{figure}

\section{Experimental Setup}
\label{sec:setup}

\subsection{Dataset}
\todo{Exact subset used (count, selection method, random seed), license, and any
preprocessing.}

\subsection{Baseline(s)}
\todo{Exact baseline(s) you compare against, and confirmation that your reproduction
matches published numbers (or an explanation of any gap).}

\subsection{Metrics \& Protocol}
\todo{Every metric used, computed how, and the exact comparison protocol (e.g.,
bpp-matching procedure for rate--distortion comparisons, train/val/test split, number
of seeds).}

\section{Results}
\label{sec:results}

\begin{table}[h]
\centering
\caption{Main quantitative comparison against the baseline.}
\label{tab:main}
\begin{tabular}{@{}lcccc@{}}
\toprule
\textbf{Method} & \textbf{Metric 1} & \textbf{Metric 2} & \textbf{Metric 3} & \textbf{Compute/Runtime} \\
\midrule
Baseline      & \todo{val} & \todo{val} & \todo{val} & \todo{val} \\
Ours          & \todo{val} & \todo{val} & \todo{val} & \todo{val} \\
\bottomrule
\end{tabular}
\end{table}

\begin{figure}[h]
\centering
% \includegraphics[width=0.48\linewidth]{figures/qual1.png}
% \includegraphics[width=0.48\linewidth]{figures/qual2.png}
\caption{\todo{Qualitative comparison. Describe exactly what is shown, including any
failure cases -- do not caption a figure with what you expected it to show.}}
\label{fig:qual}
\end{figure}

\subsection{Ablation Studies}
\label{sec:ablation}
\todo{At least one ablation isolating the effect of your main design choice. Table or
plot required, not just prose.}

\section{Discussion \& Limitations}
\label{sec:discussion}
\todo{What worked, what didn't, and why (to the best of your understanding). Be
explicit about limitations -- scope, dataset size, compute budget, anything that would
change if this were scaled up. A paper that honestly analyzes its limitations is
graded higher than one that overclaims.}

\section{Conclusion}
\label{sec:conclusion}
\todo{2--4 sentences: what you built, what you found, what you would do next with more
time/compute.}

\section*{Individual Contributions}
\label{sec:contributions}
\todo{One sentence per student stating their primary contributions, consistent with
the peer-evaluation form submitted separately.}

\bibliographystyle{unsrtnat}
\bibliography{references}

\appendix
\section{Additional Results}
\label{sec:appendix}
\todo{Optional: extra qualitative examples, full hyperparameter tables, or results
that did not fit in the main body.}

\end{document}